\documentclass[11pt]{article}
\usepackage{preamble}
\setlength{\parskip}{4pt}

\begin{document}

\daybanner{AP Statistics \textperiodcentered\ Unit 3 \textperiodcentered\ Topic 3.2 \textperiodcentered\ B: 2026-11-23 \textperiodcentered\ E: 2026-12-21 \textperiodcentered\ TEACHER KEY}{Fourteen Caption Users}

\sectionbanner{CED TETHER}

{\small
\begin{itemize}[leftmargin=1.5em,itemsep=2pt,topsep=2pt]
  \item \textbf{Skill 3.B} --- Determine parameters for probability distributions.
  \item \textbf{Skill 3.C} --- Describe probability distributions.
  \item \textbf{Skill 4.B} --- Interpret statistical calculations and findings to assign meaning or assess a claim.
  \item \textbf{EU UNC-3} --- Probabilistic reasoning allows us to anticipate patterns in data.
  \item \textbf{LO UNC-3.K} --- Determine parameters of a sampling distribution for sample proportions. [Skill 3.B]
  \item \textbf{EK UNC-3.K.1} --- For independent samples of a categorical variable from a population with proportion p, the sampling distribution of the sample proportion has mean $\mu$ = p and a standard deviation $\sigma$ = sqrt(p(1-p)/n).
  \item \textbf{EK UNC-3.K.2} --- If sampling without replacement, the standard deviation of the sample proportion is smaller than the formula above. If the sample size is less than 10\% of the population size, the difference is negligible.
  \item \textbf{LO UNC-3.L} --- Determine whether a sampling distribution for a sample proportion can be described as approximately normal. [Skill 3.C]
  \item \textbf{EK UNC-3.L.1} --- For a categorical variable, the sampling distribution of the sample proportion will have an approximate normal distribution provided the sample size is large enough: np $\geq$ 10 and n(1-p) $\geq$ 10.
  \item \textbf{LO UNC-3.M} --- Interpret probabilities and parameters for a sampling distribution for a sample proportion. [Skill 4.B]
  \item \textbf{EK UNC-3.M.1} --- Probabilities and parameters for a sampling distribution for a sample proportion should be interpreted using appropriate units and within the context of a specific population.
\end{itemize}
}
\textbf{Essential question:} How do sampling, spread, and shape conditions limit what a surprising sample proportion can show?\par
\textbf{DOK-3 skill:} 4.B \textperiodcentered\ \textbf{FRQ pattern:} {\small\texttt{failed-\allowbreak{}large-\allowbreak{}counts-\allowbreak{}normal-\allowbreak{}tail-\allowbreak{}vs-\allowbreak{}valid-\allowbreak{}binomial-\allowbreak{}and-\allowbreak{}bias}}\par
\textbf{DOK rationale:} Part (c) separates independence, shape, and unbiasedness to evaluate two overclaims and use a supplied probability to bound the conclusion supported by a random sample.\par

\frameworkphaseheader{Do Now (first take) $\rightarrow$ Explore (video + follow-along) $\rightarrow$ Exit (finish + turn in)}{1 $\rightarrow$ 3}{45}{%
\begin{itemize}[leftmargin=1.5em,itemsep=2pt,topsep=2pt]
  \item Board slide up at the bell. Ask students to mark which sentence in each claim is about shape, center, or evidence.
  \item Begin the lesson video follow-along at minute 5.
  \item Return to the board slide at minute 33. Circulate and press the distinct roles of 10 percent and Large Counts.
  \item Collect at the door. Score part (c) only, E/P/I.
\end{itemize}
}{%
\begin{itemize}[leftmargin=1.5em,itemsep=2pt,topsep=2pt]
  \item Choose whose reasoning is closer and name one condition before the lesson.
  \item Complete the follow-along about the center, spread, shape, and probabilities for sample proportions.
  \item Finish (a)--(c), revise the first take in the exit line, and turn in the sheet.
\end{itemize}
}{%
\begin{itemize}[leftmargin=1.5em,itemsep=2pt,topsep=2pt]
  \item Which condition supports treating the 100 selections as approximately independent?
  \item Why does $np=8$ address shape rather than the estimator's long-run center?
  \item What values of $X$ belong to $\widehat p\geq0.14$?
  \item How do 0.0135 and about 0.028 change the strength, but not the direction, of the evidence?
\end{itemize}
}{Solo teacher. Circulate ~5 pairs. Require students to attach each condition to its job and to distinguish evidence against 0.08 from proof of a changed population rate.}

\begin{callout}{\IconWarn\ WATCH FOR}{calloutred}
\begin{itemize}[leftmargin=1.5em,itemsep=2pt,topsep=2pt]
  \item Applying an $n>30$ rule carried over from sample means.
  \item Calling a nonnormal sampling distribution biased.
  \item Calculating exactly 14 rather than 14 or more caption users.
\end{itemize}
\end{callout}

\sectionbanner{SCORING PART (c) --- E / P / I}

\textbf{Expected elements}
\begin{itemize}[leftmargin=1.5em,itemsep=2pt,topsep=2pt]
  \item \textbf{judgment-1} (required) --- Uses the passing 10 percent condition and failed np=8 normal-shape condition to distinguish the two model requirements.
  \item \textbf{judgment-2} (required) --- Uses supplied tail near 0.0282 as unusual evidence under p=0.08, without claiming proof of an increase.
  \item \textbf{judgment-3} (required) --- Explains that SRS unbiasedness concerns center, whereas normality concerns shape, and rejects both overclaims.
\end{itemize}

{\small\renewcommand{\arraystretch}{1.25}
\noindent\begin{tabular}{|p{0.5in}|p{5.9in}|}
\hline
\textbf{E} & Meets all three checks with correct contextual reasoning; equivalent wording is acceptable. \\ \hline
\textbf{P} & Shows at least one substantive correct check, but has an omission or error that prevents E. \\ \hline
\textbf{I} & Shows none of the three checks with substantive correct reasoning; a choice alone is insufficient. \\ \hline
\end{tabular}}

\textbf{Common mistakes}
\begin{itemize}[leftmargin=1.5em,itemsep=2pt,topsep=2pt]
  \item Checking only $n>30$ instead of both expected counts
  \item Using $P(X=14)$ rather than the upper-tail event $P(X\geq14)$
  \item Saying the 10 percent condition determines normal shape rather than approximate independence and spread
  \item Equating a skewed sampling distribution with a biased estimator
\end{itemize}

\clearpage
\focusbanner{TODAY'S PROBLEM --- annotated}

Suppose a video platform has 5,000 active accounts and says that $p=0.08$ of them regularly use captions. An SRS of 100 accounts, selected without replacement, contains 14 regular caption users, so $\widehat p=0.14$.\par\smallskip \textbf{Jules:} ``Because $n=100>30$, $\widehat p$ is approximately normal. I get $P(\widehat p\geq0.14)=0.0135$, which proves caption use increased.''\par \textbf{Nora:} ``Because $np=8<10$, a normal model is not justified. Therefore $\widehat p$ is biased and this sample tells us nothing.''\par\smallskip A binomial calculator gives $P(X\geq14)=0.0282$ for $X\sim\mathrm{Bin}(100,0.08)$. This closely approximates sampling from 5,000 accounts.\par
\begin{center}
\textbf{Caption-use sample}\par\smallskip
\begin{tabular}{|l|c|c|c|}
\hline
\textbf{Source} & \textbf{Accounts} & \textbf{Users} & \textbf{Rate} \\ \hline
\textbf{Population} & 5,000 & 400 & 0.08 \\ \hline
\textbf{SRS} & 100 & 14 & 0.14 \\ \hline
\end{tabular}
\end{center}
\begin{firsttakebox}{FIRST TAKE (5 min)}
Does finding 14 caption users instead of the expected 8 prove that use increased? Give one reason.\par
\answer{Not graded. Each student has one correct observation and one overreach.}
\end{firsttakebox}

\textit{Rules callout ``CHECK INDEPENDENCE, THEN SHAPE (keep on the page)'' is printed on the student sheet.}\par\medskip

\dokbadge{1}~\textbf{(a)} Identify $p$ and verify the sample proportion $\widehat p=14/100$.\par
\answer{The parameter is $p$, the proportion of all 5,000 active accounts that regularly use captions. The statistic is the sample proportion $\widehat p=14/100=0.14$. Thus $N=5{,}000$ and $n=100$.}

\dokbadge{2}~\textbf{(b)} Find the mean and standard deviation of the sampling distribution of $\widehat p$. Check the 10 percent condition and both Large Counts values.\par
\answer{The sampling distribution has mean $\mu_{\widehat p}=p=0.08$ and standard deviation $\sigma_{\widehat p}=\sqrt{0.08(0.92)/100}=0.027129\ldots\approx0.0271$. The 10 percent condition holds because $100<0.10(5{,}000)=500$, so dependence from sampling without replacement is negligible. Large Counts does not fully hold: $np=100(0.08)=8<10$, while $n(1-p)=100(0.92)=92\geq10$.}

\dokbadge{3}~\textbf{(c)} In 3--4 sentences, explain what each student gets right and wrong. Use the condition checks and the supplied probability 0.0282. State what the sample suggests about caption use and why a failed normal-shape check does not make the estimate biased.\par
\answer{Neither claim is fully correct: the sample is random and $100<500$, but $np=8<10$, so Nora correctly rejects the normal approximation. The supplied binomial tail, about 0.0282, makes 14 or more unusual under $p=0.08$, giving evidence worth investigating rather than proof of an increase. The SRS produces an unbiased sample proportion; the failed condition concerns shape, not its long-run center. Jules overstates the evidence and Nora wrongly discards a useful estimate.}

\begin{summaryexitbox}{EXIT}
One thing the video changed about my first take: \hrulefill
\end{summaryexitbox}

\end{document}
